% ============================================================
% CAPITOLUL 6 — PREDICȚIE CU REȚELE LSTM
% ============================================================

\chapter{Predicție Epidemiologică folosind Rețele LSTM}
\label{ch:lstm}

\section{Rolul componentei LSTM}
\label{sec:rol_lstm}

Modelele compartimentale oferă interpretabilitate: fiecare compartiment are o
semnificație epidemiologică, iar parametrii descriu rate de tranziție între stări.
Această proprietate este esențială pentru o lucrare care urmărește modelele din
familia SIR și extensia SIDARTHE. Totuși, seriile epidemiologice reale conțin
tipare care nu sunt întotdeauna modelate explicit prin ecuații diferențiale:
întârzieri de raportare, efecte de weekend, modificări administrative ale
datelor, schimbări bruște de comportament sau intervenții care nu sunt codificate
direct în parametri.

Din acest motiv, o componentă LSTM este utilă ca instrument auxiliar de predicție
pe serii temporale. Ea nu înlocuiește modelele SIR/SIS/SEIR/SIRS/SIDARTHE, ci
oferă o linie de comparație statistică și o posibilă corecție a reziduurilor
rămase după rularea unui model compartimental calibrat.

\section{Motivația alegerii LSTM}
\label{sec:motivatie_lstm}

LSTM (\emph{Long Short-Term Memory}) este o arhitectură recurentă proiectată
pentru modelarea dependențelor temporale pe termen lung \cite{hochreiter1997lstm}.
Într-o serie epidemiologică, valoarea observată într-o zi depinde de un istoric
recent: numărul de cazuri din zilele anterioare, tendința de creștere sau scădere
și regularitățile de raportare. O rețea recurentă poate învăța astfel de tipare
direct din date, fără a impune explicit o structură de compartimente.

Această flexibilitate este în același timp avantaj și risc. LSTM poate aproxima
relații dificil de exprimat prin ecuații simple, dar poate învăța și artefacte
ale datelor. De aceea, în această lucrare LSTM este tratat ca instrument
predictiv experimental, nu ca model epidemiologic explicativ.

\section{Formularea problemei de predicție}
\label{sec:formulare_predictie_lstm}

Fie $y_t$ seria observată la ziua $t$, de exemplu cazuri active, cazuri noi sau
decese. O rețea LSTM primește o fereastră de lungime $\ell$ și produce o
predicție pentru un orizont $h$:
\begin{equation}
\hat{y}_{t+h} =
f_{\phi}\left(y_{t-\ell+1}, y_{t-\ell+2}, \ldots, y_t\right),
\end{equation}
unde $\phi$ reprezintă parametrii rețelei. Pentru predicție multi-pas, ieșirea
poate fi extinsă la un vector:
\begin{equation}
\left(\hat{y}_{t+1}, \hat{y}_{t+2}, \ldots, \hat{y}_{t+h}\right)
= f_{\phi}\left(y_{t-\ell+1:t}\right).
\end{equation}

În contextul aplicației, intrarea poate fi univariată sau multivariată. Varianta
univariată folosește o singură serie, de exemplu cazuri active. Varianta
multivariată poate include:
\begin{itemize}
    \item cazuri noi zilnice;
    \item cazuri active estimate;
    \item decese zilnice sau cumulative;
    \item populația regiunii;
    \item ziua săptămânii, pentru efecte de raportare;
    \item traiectorii produse de un model compartimental.
\end{itemize}

Datele trebuie împărțite cronologic. O împărțire aleatoare ar produce scurgeri de
informație, deoarece exemple din viitor ar putea ajunge în antrenare. Normalizarea
se calculează doar pe fereastra de antrenare și se aplică neschimbat pe validare
și test.

\section{Integrarea cu modelele compartimentale}
\label{sec:integrare_lstm_compartimental}

Există două moduri naturale de integrare a LSTM în sistemul propus. Primul este
modelul LSTM independent, folosit ca linie de bază predictivă. În acest caz,
rețeaua primește doar date observate și produce o predicție numerică. Rezultatul
poate fi comparat cu predicția SIR, SEIR sau SIDARTHE calibrată pe aceeași
fereastră.

Al doilea mod este modelul hibrid compartimental--LSTM. Modelul compartimental
produce o predicție interpretabilă $\hat{y}^{(c)}_t$, iar LSTM învață reziduul:
\begin{equation}
r_t = y_t - \hat{y}^{(c)}_t.
\end{equation}
Predicția finală este:
\begin{equation}
\hat{y}^{(hibrid)}_{t+h} =
\hat{y}^{(c)}_{t+h} + f_{\phi}(r_{t-\ell+1}, \ldots, r_t).
\end{equation}

Această variantă păstrează rolul central al modelului epidemiologic. LSTM nu
încearcă să explice mecanismul de transmitere, ci corectează deviațiile
statistice rămase după simularea compartimentală. Dacă reziduurile au un tipar
stabil, modelul hibrid poate îmbunătăți predicția pe termen scurt.

\section{Fluxul de antrenare și evaluare}
\label{sec:flux_lstm}

\begin{algorithm}[t]
\caption{Evaluarea unei componente LSTM pentru serii epidemiologice}
\label{alg:evaluare_lstm}
\begin{lstlisting}[style=pseudocode]
Input: dataset temporal Y, lungime fereastra l,
       orizont h, model compartimental optional m
Output: predictii LSTM, metrici MAE/RMSE/nRMSE

ordoneaza datele cronologic
imparte datele in train, validare si test
calculeaza normalizarea doar pe train
construieste exemple din ferestre de lungime l
daca exista model compartimental m:
    ruleaza modelul m si calculeaza reziduuri
    foloseste reziduurile ca tinta pentru LSTM
altfel:
    foloseste seria observata ca tinta
antreneaza LSTM pentru numarul configurat de epoci
calculeaza predictii pe fereastra de test
raporteaza MAE, RMSE si nRMSE
\end{lstlisting}
\end{algorithm}

Evaluarea trebuie făcută pe o fereastră de test care nu a fost folosită la
antrenare sau la alegerea hiperparametrilor. Pentru comparație corectă, aceeași
fereastră de test trebuie folosită pentru modelul compartimental calibrat,
pentru LSTM independent și pentru modelul hibrid.

\section{Metrici și linii de bază}
\label{sec:metrici_lstm}

Componenta LSTM folosește aceleași metrici ca analiza comparativă a modelelor:
MAE, RMSE și nRMSE. Pentru a evita concluzii exagerate, rezultatele trebuie
comparate cu linii de bază simple:
\begin{itemize}
    \item persistență: $\hat{y}_{t+1}=y_t$;
    \item medie mobilă pe ultimele $k$ zile;
    \item model compartimental calibrat;
    \item LSTM independent;
    \item model hibrid compartimental--LSTM.
\end{itemize}

O componentă LSTM este justificată doar dacă îmbunătățește fereastra de test față
de aceste repere. O eroare mică pe fereastra de antrenare nu este suficientă,
deoarece seriile epidemiologice sunt scurte și zgomotoase, iar riscul de
supraînvățare este ridicat.

\section{Locul în arhitectura aplicației}
\label{sec:arhitectura_lstm}

Integrarea practică a LSTM se face printr-un runner separat, nu prin extinderea
runner-ului Python folosit pentru modelele compartimentale. Separarea este
importantă deoarece predicția LSTM are alt profil de calcul, alți parametri,
altă logică de citire a datelor și altă semantică a rezultatelor. Runner-ul
\texttt{runner-lstm} consumă coada Redis \texttt{epimodel:jobs:lstm}, în timp ce
runner-ele R, Julia și Python rămân dedicate simulărilor ODE.

Fluxul rămâne compatibil cu arhitectura curentă:
\begin{itemize}
    \item API-ul creează un job de predicție;
    \item API-ul rutează explicit modelul \texttt{lstm} către coada
    \texttt{epimodel:jobs:lstm};
    \item Redis decuplează cererea de execuția efectivă;
    \item runner-ul LSTM citește seria \texttt{infected} din dataset;
    \item modelul LSTM este antrenat pe o fereastră glisantă de lungime
    \texttt{lookback};
    \item predicțiile și metricile sunt persistate în baza de date;
    \item frontend-ul afișează seriile \texttt{observed}, \texttt{fitted} și
    \texttt{predicted}.
\end{itemize}

Această separare păstrează clară distincția dintre simulare epidemiologică și
predicție statistică. Modelul compartimental rămâne responsabil pentru
interpretarea mecanismului, iar LSTM este folosit pentru prognoză pe termen
scurt și pentru analiza reziduurilor.

\section{Limitări}
\label{sec:limitari_lstm}

LSTM necesită mai multe date decât modelele compartimentale simple. Pentru regiuni
cu serii scurte sau raportări instabile, rețeaua poate învăța artefacte
administrative în locul dinamicii epidemiologice. De asemenea, rezultatul LSTM
este mai puțin interpretabil: o predicție bună numeric nu explică prin ce
tranziții epidemiologice se produce evoluția.

Din acest motiv, LSTM trebuie prezentat ca extensie complementară. El poate ajuta
la prognoză și la detectarea reziduurilor sistematice, dar nu înlocuiește
analiza modelelor SIR/SIS/SEIR/SIRS/SIDARTHE și nici calibrarea parametrilor
epidemiologici.

\testtodo{Implementarea rulează prin \texttt{runner-lstm}, dar pentru rezultatul
final al lucrării trebuie executată o campanie experimentală pe un dataset
concret și comparată cu persistență, medie mobilă și un model compartimental
calibrat. Rezultatele minime de raportat sunt MAE, RMSE, nRMSE, lungimea
ferestrei de intrare, orizontul de predicție, numărul de epoci și timpul de
antrenare.}
